\begin{afterword}
\summaryhead

The post asks the reader to picture a ``triangular lightbulb.'' The author reports seeing a sharp-edged pyramid, then smoothed edges, then a triangular loop of fluorescent tube, all without words. It recalls the debate over whether people really have mental images or only symbol-like labels, calls the doubting side ``spectacularly silly,'' and suggests, with ``I think,'' that it was a legacy of behaviorism. The evidence that settled the question, it says, was first mental-rotation experiments and then neuroimaging, as in Stephen Kosslyn's \textsc{Image and Brain}.

From this the post draws its thesis. The word ``apple'' is five letters, but the concept behind it can paint a picture, combine with ``cheese,'' recognize an apple and guide eating one. Words are handles for these mental paintbrushes, and through shared labels we can draw pictures in other minds. A word is a pointer, and at some point it has to be dereferenced.

\responsehead

The post's claim about imagery has held up. Mental rotation times do grow with the angle, and imagined shapes do produce shape-like activity in visual cortex. The current \textsc{Stanford Encyclopedia} entry sides with picture-like images, as the post does.

The weak part is the post's account of the debate. It names the hindsight effect and then calls the other side ``spectacularly silly,'' although it grants a few lines later that the question ``used to be harder to answer, back in the day of the controversy.'' It traces the doubting side, with ``I think,'' to behaviorism. But that side held that images are stored as symbolic internal representations, and internal representations were what behaviorism tried to keep out of psychology. The post also says behaviorism ``denied the existence of thoughts,'' which Skinner, its best-known figure, did not.

The post presents the debate as closed by 2008, on the strength of one side's book. The other side's leading figure, Pylyshyn, was still arguing in 2003 that neuroimaging did not address the format of images. Kosslyn, too, co-wrote a paper in 2015 subtitled ``Ending the imagery debate.'' The post's side may well have won since. It had not been conceded when the post declared the matter closed.

The thesis about words rests on the lightbulb, reported by introspection and hedged as such, and on the arithmetic of five bytes. It needs a limit that the post does not state. Some people report no mental imagery at all, and they use words normally. For them the title's image does not apply, while the post's own list of what a concept does, which includes recognizing and tasting, still does.

In short: a conclusion about mental imagery that has held up, reached through an account of the debate that is one-sided and judges it with hindsight.
\end{afterword}
